\section{From the tensor bound to an index inequality}
\label{sec:heat-index}

Throughout this section assume \textup{(B)}, and put
\(K=(n-1)\kappa\). Fix a bounded globally Lipschitz function
\(u\) that is constant outside a compact set. We write
\[
\Phi_t(x)=Q_tu(x)
=\inf_y\left\{u(y)+\frac{d^2(x,y)}{2t}\right\},\qquad 0<t\le1.
\]
Properness implies that the infimum is attained. Our aim is an
upper bound for the distributional Hessian of \(\Phi_1\), with
the right-hand side evaluated along the minimizing paths of this
inf-convolution.

We use the weak Hessian \eqref{eq:weak-hessian} and its weak
local Sobolev continuity for potentials in \(W^{1,2}_{\rm loc}\).
Local potentials and locally Lipschitz weights are cut off on
a neighborhood of the compact trial support when necessary.
For a local test potential \(z\), \(H_z=\Hess z\) is its
ordinary \(L^2_{\rm loc}\) Hessian.

For \(0<\delta<1\), an admissible trial is a field
\begin{equation}
\label{heat:trial-class}
Y_t=\sum_{j=1}^q a_j(t)\nabla g_j,\qquad
a_j\in C^\infty([\delta,1]),\quad
g_j\in\Test(M)\text{ compactly supported}.
\end{equation}
Here smoothness on a closed interval means restriction of a
smooth function on a neighborhood of that interval. In
particular \(Y_t=\nabla g_t\), where
\(g_t=\sum_j a_j(t)g_j\).

A geodesic \(\gamma\in\Geo(M)\) is calibrated by \(u\) if
\begin{equation}
\label{heat:calibration}
\Phi_1(\gamma_1)
=u(\gamma_0)+\tfrac12d^2(\gamma_0,\gamma_1).
\end{equation}

\begin{theorem}[Index inequality along calibrated paths]
\label{heat:index}
Let \(n\ge2\) be an integer and let \((M,d,m)\), with
\(m=\mathcal H^n\), be a full-support \(\RCD((n-1)\kappa,n)\)
space satisfying \eqref{eq:main-curvature} for every global
\(X,Y\in\TestV(M)\) and every nonnegative \(f\in\Test(M)\).
Let \(u\) be bounded globally Lipschitz
and constant outside a compact set, and let \(\Phi_t=Q_tu\).
There is a Borel prescription \(S:x\mapsto S(x)\in\Geo(M)\),
defined outside one \(m\)-null set, with \(S(x)_1=x\) and
\(S(x)\) satisfying \eqref{heat:calibration}. For every
nonnegative compactly supported Lipschitz \(f\), write
\(\pi_f=S_\#(fm)\), with forward velocity \(v\).
The plan has terminal measure \(fm\), uniformly bounded
speed, and bounded compression on every \([\delta,1]\),
\(\delta>0\).

Let \(Z\) be a compactly supported test gradient. For every
\(0<\delta<1\) and every trial \eqref{heat:trial-class}
with the global endpoint identities
\(Y_\delta=0\) and \(Y_1=Z\), one has
\begin{equation}
\label{heat:index-inequality}
H_{\Phi_1}(Z,Z)(f)
\le\int_{\Geo(M)}\int_\delta^1
\left(|\partial_tY+\nabla_vY|^2
-\kappa|Y\wedge v|^2\right)(\gamma_t)
\dd t\dd\pi_f(\gamma).
\end{equation}
The left side is the distributional evaluation
\eqref{eq:weak-hessian}; the right side is finite. Neither
the path prescription nor its velocity depends on the
choice of \(Z\), the trial, or the terminal density.
\end{theorem}

We use a forward heat potential and a backward
density factor, in the entropic-interpolation framework of
\cite{Leonard2017,GigliTamanini2021}. In particular, the density is prescribed at time
one; its equation has a backward diffusion term when expressed
using the forward potential.
Proposition~\ref{heat:viscous-index} gives the index estimate
at positive viscosity. We then identify its transport limit
and prove strong convergence of the velocities before passing
the quadratic energy to Theorem~\ref{heat:index}.

\subsection{Heat inputs and positive terminal data}

We record the heat estimates, including the extension needed on
an infinite-measure space. Constants in this subsection may depend
on \(K,n\) and a fixed upper bound for the heat time.

\begin{lemma}[Heat estimates and bounded data]
\label{heat:analytic-inputs}
Let \(P_s=\exp(s\Delta)\). On \(0<s\le S<\infty\) one has
\begin{align}
\label{heat:bakry-one}
|\nabla P_sh|&\le e^{-Ks}P_s|\nabla h|,\\
\label{heat:reverse-poincare}
|\nabla P_sh|^2&\le C_Ss^{-1}P_s(h^2),\\
\label{heat:laplacian-estimate}
|\Delta P_sh|^2&\le C_Ss^{-1}P_s(|\nabla h|^2),
\end{align}
with the usual Sobolev hypothesis for the first and third
inequalities and \(h\in L^2\) for the second. In particular,
\begin{equation}
\label{heat:laplacian-from-values}
|\Delta P_sh|^2\le C_Ss^{-2}P_s(h^2).
\end{equation}
Positive-time heat regularization of an \(L^2\cap L^\infty\)
function has bounded Lipschitz representatives for all its
generator powers. These assertions extend to bounded data, by
local approximation, whenever the right-hand sides are defined
by the heat kernel.

For every bounded nonnegative, nonzero datum \(h\), and
\(0<s\le1\), the locally defined logarithmic Laplacian satisfies
\begin{equation}
\label{heat:li-yau-consequence}
\Delta\log P_sh\ge
-C\left(s^{-1}+s|\nabla\log P_sh|^2\right).
\end{equation}
All logarithmic statements are local on regions where the
positive-time solution has a positive lower bound.
\end{lemma}

\begin{proof}
The first inequality is the \(L^1\) gradient estimate
\cite[Corollary~3.5]{Savare2014}; the second is the reverse
Poincar\'e inequality. The Bakry--Ledoux estimate
\cite[Proposition~4.9]{ErbarKuwadaSturm2015} gives the third:
its coefficient of \(|\Delta P_sh|^2\), before rearrangement,
is \((1-e^{-2Ks})/(nK)\), with value \(2s/n\) at \(K=0\).
Apply \eqref{heat:laplacian-estimate} at time \(s/2\) to
\(P_{s/2}h\), then apply \eqref{heat:reverse-poincare} inside
the remaining heat operator, to obtain
\eqref{heat:laplacian-from-values}.

These inequalities give the \(L^\infty\)-to-Lipschitz and
\(L^\infty\)-to-\(L^\infty\) Laplacian bounds at positive time.
Semigroup splitting, together with commutation in \(L^2\),
therefore gives, for each integer \(j\ge0\),
\[
\|\Delta^jP_sh\|_\infty
+\sqrt{s}\,\Lip(\Delta^jP_sh)
\le C_{j,S}s^{-j}\|h\|_\infty.
\]
The time dependence is differentiable in the corresponding
\(L^2\) generator domains on intervals away from zero.

Here is a local extension argument that does not assume
\(m(M)<\infty\). Choose \(0\le\chi_R\le1\), equal to one on
\(B_R(o)\), compactly supported, and set \(h_R=\chi_Rh\).
For \(R'\ge R\), the difference \(d_R=h_R-h_{R'}\) vanishes
on \(B_R(o)\) and satisfies
\(|d_R|^2\le\|h\|_\infty^2\mathbf1_{M\setminus B_R(o)}\).
The Gaussian kernel bound shows that \(P_s(d_R^2)\) tends to
zero uniformly on compact spatial sets and compact subintervals
of \(s>0\). Equations
\eqref{heat:reverse-poincare} and
\eqref{heat:laplacian-from-values} then make the gradients and
Laplacians Cauchy locally in \(L^\infty\), as well as the
functions themselves. Distributional closure identifies their
limits with those of \(P_sh\). Repeated splitting gives the
same local conclusions for generator powers. The kernel bounds
used here are valid for all centers, with constants depending
only on \(K,n,S\); see
\cite[Theorem~1.2]{JiangLiZhang2016}.

For integrable nonnegative data, the general-measure Li--Yau
inequality is \cite[Theorem~3.4]{GigliTamanini2021}:
\[
|\nabla\log P_sh|^2
\le e^{-2Ks/3}\frac{\Delta P_sh}{P_sh}
+\frac{nK}{3}\frac{e^{-4Ks/3}}{1-e^{-2Ks/3}}.
\]
Its value at \(K=0\) is understood by continuity. We may weaken
the Ricci bound to a strictly negative constant. Rearranging
the inequality, using \(|e^{2Ks/3}-1|\le Cs\) and the
\(C/s\) bound for the remaining coefficient on \(s\le1\),
gives \eqref{heat:li-yau-consequence}. Apply this to \(h_R\)
and pass to the local limit just established. Positive-time
positivity and continuity supply a lower bound on each compact
region, so division and the logarithmic chain rule pass to the
limit. This proves the assertion also for bounded data that are
not in any finite \(L^p\) space.
\end{proof}
Fix \(o\in M\) and choose \(D>0\) sufficiently large that
\begin{equation}
\label{heat:reference-density}
\bar f(x)=e^{-Dd(o,x)},\qquad
\int(1+d^2(o,x))\bar f(x)\dd m(x)<\infty.
\end{equation}
Such a choice follows from Bishop--Gromov volume growth.
The initial positive terminal class will be
\begin{equation}
\label{heat:positive-data}
\mathcal F=
\{f_0+\lambda\bar f:
  f_0\ge0\text{ compactly supported Lipschitz},\ \lambda>0\}.
\end{equation}
Every member is bounded, belongs to \(L^1\cap W^{1,2}\), has
finite second moment, and has a globally Lipschitz logarithm.
For the last assertion, on a compact neighborhood of
\(\supp f_0\) the denominator \(f\) has a positive lower bound,
and outside it \(\log f=\log\lambda-Dd(o,\cdot)\).
The resulting uniform local Lipschitz bound is global because
the space is a length space. The class is closed under positive
rescaling. We normalize \(\int f\dd m=1\) when using
Wasserstein distances; all final inequalities are homogeneous
in the terminal mass.

For \(f\in\mathcal F\) and \(0<c\le1\), define
\begin{align}
\label{heat:two-factors}
H_t^c&=P_{ct}(e^{-u/(2c)}),&
J_t^c&=P_{c(1-t)}(f/H_1^c),&
\rho_t^c&=H_t^cJ_t^c,\\
\label{heat:potentials}
\phi^c&=-2c\log H^c,&
\psi^c&=2c\log J^c,&
b^c&=(\psi^c-\phi^c)/2=c\log\rho^c,\\
\theta^c&=(\psi^c+\phi^c)/2,&
v_c&=\nabla\theta^c,&w_c&=\nabla b^c.
\end{align}
Here \(H_t^c\) is a scalar heat factor; the notation \(H_z\)
continues to denote the Hessian of a potential \(z\). The field
\(v_c\) is the continuity-equation velocity, while
\(w_c=c\nabla\log\rho^c\) records the density gradient. We suppress
the superscript \(c\) in calculations at fixed viscosity.

\begin{lemma}[Regularity, equations, and endpoint traces]
\label{heat:regularity}
For each fixed \(c>0\), on compact spatial sets and compact
time intervals in \((0,1)\), the factors in
\eqref{heat:two-factors} have bounded Lipschitz generator powers
and positive lower bounds. Their logarithms are local test
functions, with locally bounded Sobolev Laplacians. The density
is a positive local test function. Time differentiation of the
potentials is valid locally in the Laplacian graph norm, hence
also for their \(L^2\) Hessians.

There is a constant \(L\), independent of \(c\) and \(t\), such
that
\begin{equation}
\label{heat:uniform-gradients}
|\nabla\phi|+|\nabla\psi|+|v_c|+|w_c|\le L.
\end{equation}
The following identities hold on \((0,1)\):
\begin{align}
\label{heat:continuity}
\partial_t\rho&=-\divg(\rho v_c),&
\nabla\rho&=\rho w_c/c,\\
\label{heat:hamilton-jacobi}
\partial_t\theta&=-\tfrac12(|v_c|^2+|w_c|^2)-c\Delta b,&
\partial_t\phi&=-\tfrac12|\nabla\phi|^2+c\Delta\phi.
\end{align}
Moreover, \(\int\rho_t\dd m=1\), \(\phi_t\) is uniformly
bounded, and, for fixed \(c\), as \(t\uparrow1\),
\begin{equation}
\label{heat:terminal-traces}
\rho_t\longrightarrow f,\qquad
\theta_t\longrightarrow\phi_1+c\log f
\quad\text{locally uniformly and strongly in }W^{1,2}_{\rm loc}.
\end{equation}
The local Lipschitz bounds for \(\rho_t\) near the endpoint
may depend on the fixed \(c\). The bound for
\(\nabla\theta_t\) is uniform for \(0<c\le1\), and the
terminal density is always the fixed function \(f\).
\end{lemma}

\begin{proof}
The first heat datum is a constant plus a compactly supported
\(W^{1,2}\cap L^\infty\) function. Constants are preserved by
the heat flow, so Lemma~\ref{heat:analytic-inputs} applies.
Since \(u\) is bounded, \(H_t\) has positive global lower and
upper bounds for fixed \(c\). Thus \(g=f/H_1\) belongs to
\(L^1\cap W^{1,2}\cap L^\infty\), and the same lemma applies
to \(J\). Positivity of the kernel gives \(J>0\); its continuous
representative has a positive minimum on each compact
space-time region under consideration.

The logarithmic Laplacian formula
\[
\Delta\log H=\frac{\Delta H}{H}
-\frac{|\nabla H|^2}{H^2}
\]
and its counterpart for \(J\) give local boundedness and
\(W^{1,2}\) regularity. The squared-gradient term is Sobolev
by metric compatibility and the \(L^2\) Hessian estimate.
Product and chain rules similarly show that
\(\rho\) is a local test function. Generator differentiation
before composition, followed by these same rules, gives strong
local differentiation in the Laplacian graph norm; the Hessian
estimate \eqref{eq:pre-hessian-estimate} then gives the Hessian
statement.

If \(L_u=\Lip(u)\), then
\(|\nabla e^{-u/(2c)}|\le (L_u/(2c))e^{-u/(2c)}\).
Use \eqref{heat:bakry-one}, applied to the compact perturbation
of the constant, to get
\(|\nabla\phi_t|\le e^{-Kct}L_u\).
Also
\[
\Lip(\log g)\le\Lip(\log f)+C/c.
\]
Applying \eqref{heat:bakry-one} to \(g\) gives the corresponding
bound for \(\nabla\psi\). This proves
\eqref{heat:uniform-gradients} and also bounds the gradients of
the terminal expressions. Since \(\rho=HJ\),
\(\nabla\rho=\rho w_c/c\); differentiating the two factors
proves \eqref{heat:continuity}--\eqref{heat:hamilton-jacobi}
with the displayed signs.

Symmetry and the semigroup property, with nonnegative
integrands, give
\[
\int H_tP_{c(1-t)}(f/H_1)\dd m
=\int H_1(f/H_1)\dd m=1.
\]
The maximum principle bounds \(\phi_t\) between \(\inf u\)
and \(\sup u\). Finally \(J_t\to g\) strongly in \(W^{1,2}\).
Heat continuity for bounded Lipschitz data, using the uniform
small-time first moment of the kernel, gives uniform
convergence of \(J_t\) to \(g\). Strong local convergence for
the products and logarithms follows from their local positive
lower bounds. This proves \eqref{heat:terminal-traces}.
\end{proof}

\Needspace{6\baselineskip}
\subsection{Uniform density and weighted second-order bounds}

Two different controls are needed below. Bounded densities allow
pairing with fixed Sobolev coefficients and passage to the transport
limit. A one-sided bound for \(c\Delta b\) will instead control its
product with a nonnegative Bochner measure in the Riccati calculation.
That bound must remain integrable up to the prescribed terminal time.

\begin{lemma}[The terminal-time kernel ratio]
\label{heat:kernel-ratio}
For each \(A<\infty\) there is \(C_A<\infty\) such that, if
\(0<s\le1\), \(g>0\) is bounded and
\(\Lip(\log g)\sqrt{s}\le A\), then
\begin{equation}
\label{heat:kernel-ratio-estimate}
\frac{(P_sg^2)^{1/2}}{P_sg}\le C_A
\qquad\text{everywhere.}
\end{equation}
The constant is independent of the base point and of a positive
lower bound for \(g\).
\end{lemma}

\begin{proof}
The two-sided Gaussian estimate and Bishop--Gromov give constants
\(a,C,c_0>0\), uniform in the base point, for which
\[
p_s(x,y)\le\frac{C}{m(B_{\sqrt{s}}(x))}
e^{-d^2(x,y)/(as)},\qquad
p_s(x,y)\ge\frac{c_0}{m(B_{\sqrt{s}}(x))}
\quad(y\in B_{\sqrt{s}}(x)).
\]
The harmless bounded factors involving \(s\le1\) are absorbed
in the constants. In addition, for \(j\ge0\),
\[
\frac{m(B_{(j+1)\sqrt{s}}(x))}{m(B_{\sqrt{s}}(x))}
\le C(j+1)^n e^{C(j+1)\sqrt{s}}.
\]
The lower bound and the logarithmic Lipschitz condition give
\(P_sg(x)\ge c_0e^{-A}g(x)\). On the annulus
\(j\sqrt{s}\le d(x,y)<(j+1)\sqrt{s}\),
\(g(y)^2\le g(x)^2e^{2A(j+1)}\). Therefore
\[
P_sg^2(x)\le Cg(x)^2
\sum_{j=0}^\infty(j+1)^n
\exp\bigl(-j^2/a+(C+2A)(j+1)\bigr).
\]
This convergent series is uniform in \(x,s\). Dividing the two
estimates proves the claim.
\end{proof}

\begin{lemma}[Density and lower Laplacian control]
\label{heat:density-laplacian}
For each \(0<\delta<1\) there are \(C_\delta\) and
nonnegative functions \(e_c\) on \((\delta,1)\) such that
\begin{equation}
\label{heat:density-and-error}
\sup_{\delta\le t\le1}\|\rho_t^c\|_\infty\le C_\delta,
\qquad c\Delta b^c\ge-e_c(t),
\qquad \|e_c\|_{L^1(\delta,1)}\longrightarrow0.
\end{equation}
The lower bound holds also for quasi-continuous representatives,
for almost every time.
\end{lemma}

\begin{proof}
\emph{Bounded compression.}
Apply \eqref{heat:li-yau-consequence} at heat time \(ct\),
using \eqref{heat:uniform-gradients}. It gives
\[
\Delta\phi_t\le C(1+t^{-1}),\qquad
-\Delta\psi_t\le C(1+(1-t)^{-1}).
\]
The constants are independent of \(c\). Since
\(v_c=\nabla\phi+w_c\), the density equation can also be written
\begin{equation}
\label{heat:backward-fokker-planck}
\partial_t\rho=-\divg(\rho\nabla\phi)-c\Delta\rho.
\end{equation}
For \(p\ge2\), multiplication by \(p\rho^{p-1}\) gives
\[
\frac{d}{dt}\int\rho^p\dd m
=-(p-1)\int\rho^p\Delta\phi\dd m
+cp(p-1)\int\rho^{p-2}|\nabla\rho|^2\dd m
\ge-(p-1)C(1+t^{-1})\int\rho^p\dd m.
\]
For fixed \(c\) the integrations are justified first with
exhausting cutoffs: \(\rho\) is bounded and integrable,
\(|\nabla\rho|\le C_c\rho\), and \(\Delta\phi\) is globally
bounded away from the initial time. The cutoff errors vanish
by integrability. One may equivalently use the strong interior
\(L^2\) heat differentiation before removing the cutoffs.
Integrate backward from \(\rho_1=f\), then let \(p\to\infty\).
This proves the density bound, with the direction of integration
fixed by \eqref{heat:backward-fokker-planck}.

\emph{The terminal lower bound.}
For the endpoint bound let \(r=1-t\), \(g=f/H_1\), and
\(s=cr\). When \(r\le c\), the logarithmic Lipschitz bound for
\(g\) is at most \(C/c\), and \(s\le c^2\). Lemma
\ref{heat:kernel-ratio} and
\eqref{heat:laplacian-estimate} consequently give
\[
\frac{|\Delta P_sg|}{P_sg}
\le C s^{-1/2}c^{-1},\qquad
c|\Delta\psi_t|
\le C\left(1+\sqrt{c/r}\right).
\]
For \(r\ge c\), use the preceding Li--Yau bounds instead.
Since \(b=(\psi-\phi)/2\), one can thus take, after increasing
constants depending on \(\delta\),
\[
e_c(t)\le
\begin{cases}
C_\delta c(1+r^{-1}),&r\ge c,\\
C_\delta(c+1+\sqrt{c/r}),&0<r<c.
\end{cases}
\]
For sufficiently small \(c\), its integral is at most
\(C_\delta(c|\log c|+c)\). The Sobolev function \(\Delta b\)
has the same lower bound outside a polar set in its
quasi-continuous representative, by the capacitary facts in
Section~\ref{sec:preliminaries}.
\end{proof}

The next estimates adapt the weighted Hessian and Laplacian bounds,
and the vanishing density-gradient estimates, of
\cite[Lemmas~4.9--4.10]{GigliTamanini2021}. Those estimates are stated
away from both endpoints. Here the fixed terminal density and its
first-order trace let us integrate all the way to time one; the
argument below includes that endpoint step.

\begin{lemma}[Weighted second-order estimates]
\label{heat:weighted-estimates}
For every bounded set \(E\) and \(0<\delta<1\),
\begin{equation}
\label{heat:weighted-estimate}
\int_\delta^1\!\int_E\rho^c
\left(|H_{b^c}|^2+|H_{\theta^c}|^2
+|\Delta b^c|^2+|\Delta\theta^c|^2
+\frac{|w_c|^4}{c^2}\right)\dd m\dd t
\le C_{\delta,E}.
\end{equation}
Furthermore, for every fixed \(T\in L^2((\delta,1)\times E)\),
\begin{equation}
\label{heat:weighted-small-w}
\int_\delta^1\!\int_E\rho^c|w_c|^2|T|^2\dd m\dd t
\longrightarrow0.
\end{equation}
\end{lemma}

\begin{proof}
\emph{The integrated Bochner identity.}
Choose a nonnegative space-time cutoff \(\chi\), zero for
\(t\le\delta/2\), equal to one on \([\delta,1]\times E\),
and compactly supported in space, such that
\[
|\partial_t\chi|+|\nabla\chi|+|\Delta\chi|
\le C\sqrt\chi.
\]
For example take fourth powers of a smooth time cutoff and a
good spatial test cutoff. We first work on an interval with
terminal time \(T<1\). The equations and the definition of
\(\Gamma_2\) yield
\begin{align}
\label{heat:weighted-bochner-identity}
\frac d{dt}\int\rho\chi\Delta\theta\dd m
={}&\int\rho(\partial_t\chi+
\langle\nabla\chi,v_c\rangle)\Delta\theta\dd m\notag\\
&-\bigl(\Gamma_2(\theta)+\Gamma_2(b)\bigr)(\rho\chi)
-\int\rho\Delta b\bigl(c\Delta\chi+
\langle\nabla\chi,w_c\rangle\bigr)\dd m.
\end{align}
To check the last term, put \(q=\Delta b\). Twice integrating
by parts, using \(\nabla\rho=\rho w_c/c\), gives
\[
\int\rho\chi(\langle w_c,\nabla q\rangle+c\Delta q)
=\int\rho q(\langle w_c,\nabla\chi\rangle+c\Delta\chi).
\]
This identity is distributional; \(q\) is locally Sobolev and
bounded, and the required multipliers are admitted by the
Bochner measure calculus. Thus
\eqref{heat:weighted-bochner-identity} does not assume a
pointwise fourth derivative.

\emph{Weighted second-order bounds.}
Average \eqref{eq:bochner-hessian} and
\eqref{eq:bochner-dimension}. For either \(z=b,\theta\),
\[
\Gamma_2(z)\ge
\left(K|\nabla z|^2+\tfrac12|H_z|^2+
\tfrac1{2n}|\Delta z|^2\right)m.
\]
The two cutoff terms in
\eqref{heat:weighted-bochner-identity} are bounded by
\(\eps\int\rho\chi(|\Delta b|^2+|\Delta\theta|^2)\)
plus a constant, using \eqref{heat:uniform-gradients},
\eqref{heat:density-and-error}, and Young's inequality.
The possibly negative \(K\)-terms have the same harmless
bound. At the initial endpoint the primitive is zero. At the
terminal endpoint its limit is
\[
-\int\nabla(f\chi_1)\cdot
\bigl(\nabla\phi_1+c\nabla\log f\bigr)\dd m,
\]
which is bounded uniformly in \(c\). This follows from
\eqref{heat:terminal-traces}; no Hessian of \(\log f\) is
needed. After absorption, let \(T\uparrow1\). Monotone
convergence for the squared terms gives the first four bounds
in \eqref{heat:weighted-estimate}. In particular these bounds
also hold uniformly for all shorter terminal intervals.

\emph{The fourth-order bound.}
Choose a spatial cutoff \(\eta\) equal to one on \(E\), with
\(|\nabla\eta|\le C\sqrt\eta\). The divergence of
\(\rho\eta|w_c|^2w_c\) gives
\begin{align*}
\frac1c\int\rho\eta|w_c|^4
={}&-\int\rho\eta\bigl(|w_c|^2\Delta b
+2H_b(w_c,w_c)\bigr)\\
&-\int\rho|w_c|^2\langle w_c,\nabla\eta\rangle.
\end{align*}
All terms are integrable on interior intervals by local test
calculus. Integrate in time and then exhaust the terminal
interval. If \(A_c=\int_\delta^1\int\rho\eta|w_c|^4\),
Cauchy--Schwarz and the bounds just proved show
\(A_c/c\le C\sqrt{A_c}\). For the cutoff term use also
the uniform bound on \(w_c\) and the mass bound. Hence
\(A_c\le C^2c^2\), proving the last estimate.

\emph{Vanishing weighted products.}
It follows that \(\int_\delta^1\int_E\rho^c|w_c|^2\to0\).
For \eqref{heat:weighted-small-w}, truncate the fixed function
\(T\): on \(\{|T|\le A\}\) use this convergence, and on
its complement use the uniform bounds on \(\rho^c,w_c\).
The latter contribution is bounded by
\(C\int_{\{|T|>A\}}|T|^2\), which tends to zero as
\(A\to\infty\).
\end{proof}

\subsection{A distributional Riccati calculation}

Fix \(0<\delta<1\) and an admissible trial
\eqref{heat:trial-class}. The field and its time derivatives
are bounded and have a common compact support, while
\(\nabla Y_t,\divg Y_t,\nabla\divg Y_t\) have the stated
\(L^2\) regularity, uniformly in time. A bound on
\(\divg Y_t\) in \(L^\infty\) is not assumed.

Whenever a local heat potential or the density occurs in a
curvature evaluation with \(Y\), it is multiplied by test
cutoffs unchanged on a neighborhood of the common support of
\(Y\). Lemma~\ref{heat:regularity} and
Remark~\ref{pre:bare-gradients} then make all fields and
nonnegative multipliers admissible in \textup{(B)}. Locality
and tensoriality make the result independent of the cutoffs.

\begin{lemma}[Weak differentiated-Hessian identity]
\label{heat:riccati-identity}
Let \(z\) be a local test function, \(a=\nabla z\), \(Y\) a
compact test gradient, and \(h\) a compact Lipschitz
multiplier. Define
\begin{equation}
\label{heat:covariant-hessian-distribution}
(\nabla_aH_z)(Y,Y)(h)
=-\int H_z(Y,Y)\divg(ha)\dd m
-2\int hH_z(\nabla_aY,Y)\dd m,
\end{equation}
whenever these products are integrable. Then
\begin{equation}
\label{heat:riccati-algebra}
H_{|a|^2/2}(Y,Y)(h)
=(\nabla_aH_z)(Y,Y)(h)
+\int h|H_zY|^2\dd m+R(Y,a,a,Y)(h).
\end{equation}
For the heat potentials and density of
Lemma~\ref{heat:regularity}, every admissible trial satisfies
on \((\delta,1)\)
\begin{align}
\label{heat:riccati-evolution}
\frac d{dt}\int\rho H_\theta(Y,Y)\dd m
={}&\int\rho\left[
2H_\theta(\partial_tY+\nabla_{v_c}Y,Y)-|H_\theta Y|^2
\right]\dd m\notag\\
&-R(Y,v_c,v_c,Y)(\rho)
-H_{|w_c|^2/2+c\Delta b}(Y,Y)(\rho).
\end{align}
\end{lemma}

\begin{proof}
All assertions are local, so first insert the cutoffs just
described. Since \(\nabla_a a=H_za=\nabla(|a|^2/2)\), the
first two terms of the defining expression for
\(R(Y,a,a,Y)(h)\) sum to
\(H_{|a|^2/2}(Y,Y)(h)\). Its other terms sum to
\[
\int\divg(ha)H_z(Y,Y)
+2hH_z(\nabla_aY,Y)-h|H_zY|^2\dd m,
\]
because \([Y,a]=H_zY-\nabla_aY\). This is exactly
\eqref{heat:riccati-algebra}. The calculation uses Hessian
symmetry and first-order products only; it does not require a
third covariant derivative of \(z\).

For fixed \(c\), on a compact interior time interval the
pairing on the left of \eqref{heat:riccati-evolution} is
absolutely continuous. This follows from local graph-norm
differentiation of \(\theta\), the Hessian estimate, and local
boundedness of \(\partial_t\rho\). Differentiate that pairing,
use \(\partial_t\rho=-\divg(\rho v_c)\), and interpret its
contraction against \(H_\theta\) through
\eqref{heat:covariant-hessian-distribution}. Thus no continuity
equation is being tested against an unproved Sobolev
representative of \(H_\theta(Y,Y)\). Substitute
\eqref{heat:hamilton-jacobi} and apply
\eqref{heat:riccati-algebra} to \(z=\theta\). The
\(\nabla_{v_c}H_\theta\) terms cancel, leaving
\eqref{heat:riccati-evolution}.
\end{proof}

Completing the square in \eqref{heat:riccati-evolution} and applying
the tensor lower bound give the desired index integrand, except for
the last Hessian term. We need a uniform lower bound tending to
zero for its time integral. Smallness of \(w_c\) alone does not suffice:
differentiating \(\rho\) introduces \(1/c\). The next calculation
cancels these terms before estimating the remaining Sobolev products;
its singular measure term is controlled by the one-sided bound for
\(c\Delta b\).

\begin{lemma}[A uniform lower bound for the acceleration error]
\label{heat:error-cancellation}
For every admissible trial \(Y\), there are numbers
\(\varepsilon_c\to0\) such that, for every
\(\delta<T<1\),
\begin{equation}
\label{heat:error-lower-bound}
\int_\delta^T
H_{|w_c|^2/2+c\Delta b^c}(Y_t,Y_t)(\rho_t^c)\dd t
\ge-\varepsilon_c.
\end{equation}
The bound is independent of the terminal exhaustion parameter
\(T\).
\end{lemma}

\begin{proof}
\emph{Exact cancellations.}
We work first at fixed \(c\) and at an interior time, suppressing
these parameters. By \eqref{heat:riccati-algebra} with \(z=b\),
the error is the sum of
\[
\int\rho|H_bY|^2\dd m,\qquad
R(Y,w,w,Y)(\rho),\qquad
I=(\nabla_wH_b)(Y,Y)(\rho)+cH_{\Delta b}(Y,Y)(\rho).
\]
The first term is nonnegative. The curvature bound gives
\(R(Y,w,w,Y)(\rho)\ge
\kappa\int\rho|Y\wedge w|^2\dd m\), whose possible
negative contribution tends to zero after time integration by
\eqref{heat:weighted-small-w}.

We examine \(I\). Formula
\eqref{heat:covariant-hessian-distribution} and
\(\nabla\rho=\rho w/c\) give
\begin{equation}
\label{heat:error-transport}
(\nabla_wH_b)(Y,Y)(\rho)
=-\int\rho\left[
2H_b(\nabla_wY,Y)+H_b(Y,Y)
\left(\Delta b+\frac{|w|^2}{c}\right)\right]\dd m.
\end{equation}
Put \(D=\divg Y\), \(A=\nabla_YY\), and
\(\alpha=Yb=\langle Y,w\rangle\). Twice using the
definition of the distributional Hessian gives
\begin{align}
\label{heat:error-double-divergence}
cH_{\Delta b}(Y,Y)(\rho)
={}&\int c\Delta b\,
\left[\divg\bigl(Y\divg(\rho Y)\bigr)
+\divg(\rho A)\right]\notag\\
={}&\int\rho\Delta b
\left[H_b(Y,Y)+\frac{\alpha^2}{c}
+2\langle w,A\rangle+2\alpha D
+c\bigl(D^2+YD+\divg A\bigr)\right].
\end{align}
The last divergence is a measure, not necessarily a function.
Its product with \(\rho c\Delta b\) is understood using
quasi-continuous representatives.

To justify \eqref{heat:error-double-divergence} for the actual
test class, observe that the function
\(\divg(\rho Y)=\rho D+\rho\alpha/c\) belongs locally to
\(W^{1,2}\). The product rule for its product with the bounded
field \(Y\), proved first after truncation, gives an \(L^1\)
divergence. The multiplier \(\Delta b\) is locally bounded
Sobolev, so the truncations pass by \(L^2\) pairing for the
gradient term and dominated convergence for the divergence
term. Also \(A=\nabla(|Y|^2/2)\), whose divergence is the
Bochner Laplacian measure; multiplication by the Lipschitz
function \(\rho\) obeys its usual weak product rule.
Expanding \(Y\alpha=H_b(Y,Y)+\langle w,A\rangle\) gives
the second line of \eqref{heat:error-double-divergence}.

The terms \(\rho\Delta b H_b(Y,Y)\) in
\eqref{heat:error-transport} and
\eqref{heat:error-double-divergence} cancel. There is a second
exact cancellation:
\begin{align}
\label{heat:inverse-c-cancellation}
\int\rho\bigl(\alpha^2\Delta b-|w|^2H_b(Y,Y)\bigr)\dd m
=\int\rho\bigl[
-2\alpha\langle w,\nabla_wY\rangle
+|w|^2\alpha D+|w|^2\langle w,A\rangle\bigr]\dd m.
\end{align}
Indeed integration by parts in the first term on the left
produces
\(-\rho\alpha^2|w|^2/c-2\rho\alpha H_b(w,Y)
-2\rho\alpha\langle w,\nabla_wY\rangle\).
Writing \(H_b(Y,Y)=Y\alpha-\langle w,A\rangle\) and
integrating the negative Hessian term gives the opposites of
the first two expressions, which cancel them, with the
remaining terms displayed on the right
of \eqref{heat:inverse-c-cancellation}.

\emph{The remaining Sobolev terms.}
After division by \(c\), each term on that right side is
bounded in absolute value, on the whole time interval, by a
constant times
\[
\left(\int_\delta^1\!\int_E
\rho\frac{|w|^4}{c^2}\right)^{1/2}
\left(\int_\delta^1\!\int_E
\rho|w|^2|T_0|^2\right)^{1/2},
\]
where \(E\) contains the common support and \(T_0\) is a fixed
\(L^2\) coefficient formed from \(\nabla Y,D\) and bounded
trial fields. This tends to zero by Lemma
\ref{heat:weighted-estimates}. The remaining terms involving
\(H_b\) or \(\Delta b\), one factor \(w\), and derivatives
of \(Y\) are bounded the same way, using their weighted
\(L^2\) bound in place of the first factor.

\emph{The Bochner measure term.}
It remains to control the final parenthesis in
\eqref{heat:error-double-divergence}. Here \(D\) is only \(L^2\),
so \(D^2\) need not be an \(L^2\) coefficient. We retain it in a
nonnegative measure, to which the lower bound for \(c\Delta b\)
applies, and leave an \(L^2\) remainder. Since \(Y=\nabla g_t\),
the required decomposition is
\begin{equation}
\label{heat:positive-measure-split}
D^2+YD+\divg A=\mathcal M_t+\ell_t m,
\quad
\begin{cases}
\mathcal M_t=(\Delta g_t)^2m+\Gamma_2(g_t)-K|Y_t|^2m\ge0,\\
\ell_t=2\langle Y_t,\nabla\Delta g_t\rangle+K|Y_t|^2.
\end{cases}
\end{equation}
Here \(\ell_t\) has a uniformly bounded local \(L^2\) norm.
The measures \(\mathcal M_t\) have uniformly bounded mass
on \(E\): polarize \(\Gamma_2\) on the fixed finite list of
test potentials and use boundedness of the smooth temporal
coefficients. Equivalently, test their defining formula with
a fixed outer cutoff. Thus no \(L^4\) bound on \(D\) is used.

\emph{Time measurability and terminal uniformity.}
To justify time integration of the singular part, fix a compact
neighborhood \(E_*\) of the trial support and define the finite
signed coefficient measures and their common dominating measure by
\[
\mathcal M_{ij}=(\Delta g_i)(\Delta g_j)m
+\Gamma_2(g_i,g_j)-K\langle\nabla g_i,\nabla g_j\rangle m,
\qquad
\Lambda=\sum_{i,j}|\mathcal M_{ij}|\big|_{E_*}.
\]
Every coefficient measure neglects polar sets. The identity
\(\mathcal M_t=\sum_{i,j}a_i(t)a_j(t)\mathcal M_{ij}\)
shows that \(\mathcal M_t|_{E_*}\) has a jointly measurable,
uniformly bounded \(\Lambda\)-density.

Choose a compact test cutoff \(\chi=1\) near \(E_*\).
At fixed \(c\), on each compact interior time interval,
\(\chi\rho_t c\Delta b_t\) is strongly continuous in
\(W^{1,2}\) and has a common \(L^\infty\) bound. Indeed,
for either positive heat factor \(U\), the spatial formula
\[
\nabla\Delta\log U
=\frac{\nabla\Delta U}{U}
-\frac{\Delta U\nabla U}{U^2}
-\frac{2H_U\nabla U}{U^2}
+\frac{2|\nabla U|^2\nabla U}{U^3}
\]
expresses the derivative using bounded Lipschitz heat powers,
local positive lower bounds, and the strongly locally
\(L^2\)-continuous Hessian \(H_U\). The latter continuity
follows from the Laplacian graph estimate; continuity of the
other factors follows by integrating the next generator
powers. The fixed cutoff and \(\rho_t\) preserve these
Sobolev continuity and boundedness properties.

Strong Sobolev convergence with a common bound has
quasi-everywhere convergent subsequences. Since \(\Lambda\)
is finite and neglects polar sets, dominated convergence and
the subsequence criterion give continuity of these
quasi-continuous multipliers in \(L^1(\Lambda)\). Choose a
jointly measurable representative of this curve. At almost
every time it agrees \(\Lambda\)-almost everywhere with the
quasi-continuous multiplier. The lower bound for
\(c\Delta b_t\) therefore holds \(\Lambda\)-almost everywhere
at each admissible time, and Fubini applies without an
uncountable union of polar exceptional sets. All the bounds
in this paragraph are only fixed-\(c\), interior-time bounds.

The quasi-continuous lower bound in
\eqref{heat:density-and-error}, and the density bound there,
give
\[
\int\rho_t c\Delta b_t\dd\mathcal M_t
\ge-C_\delta e_c(t)\mathcal M_t(E).
\]
Bochner measures neglect polar sets, so an almost-everywhere
bound on the Sobolev multiplier has been used here only after
passing to its quasi-continuous representative. In particular,
for every \(T<1\),
\[
\int_\delta^T\!\int\rho_t c\Delta b_t\dd\mathcal M_t\dd t
\ge-C_\delta\sup_t\mathcal M_t(E)
\|e_c\|_{L^1(\delta,1)}.
\]
This lower error bound tends to zero uniformly in \(T\).
The absolutely continuous remainder is bounded by
\[
c\left(\int_\delta^1\!\int_E\rho|\Delta b|^2\right)^{1/2}
\left(\int_\delta^1\!\int_E\rho|\ell_t|^2\right)^{1/2}
\longrightarrow0.
\]
The Sobolev remainders were estimated in absolute value on
\((\delta,1)\), and the remaining measure contribution was bounded
below using \(\|e_c\|_{L^1(\delta,1)}\). After discarding the
nonnegative terms, these bounds hold with the same error for every
\(T<1\), proving \eqref{heat:error-lower-bound}.
\end{proof}

\begin{proposition}[The viscous index estimate]
\label{heat:viscous-index}
If the trial \eqref{heat:trial-class} satisfies \(Y_\delta=0\),
then
\begin{align}
\label{heat:viscous-index-estimate}
H_{\phi_1^c+c\log f}(Y_1,Y_1)(f)
\le\int_\delta^1\!\int\rho_t^c
\left(|\partial_tY+\nabla_{v_c}Y|^2
-\kappa|Y\wedge v_c|^2\right)\dd m\dd t+\varepsilon_c,
\end{align}
where \(\varepsilon_c\to0\).
\end{proposition}

\begin{proof}
Integrate \eqref{heat:riccati-evolution} from \(\delta\) to
\(T<1\). The initial pairing vanishes. Apply \textup{(B)}
to its curvature term and use
\[
2H_\theta(D_cY,Y)-|H_\theta Y|^2\le|D_cY|^2,
\qquad D_cY=\partial_tY+\nabla_{v_c}Y.
\]
Lemma~\ref{heat:error-cancellation} bounds the last term
uniformly in \(T\). The resulting right-hand side is
integrable up to one, by boundedness of \(\rho,v_c\) and the
fixed \(L^2\) derivatives of \(Y\).

For the terminal left-hand side, use the first-order expression
\eqref{eq:weak-hessian}. The convergences
\eqref{heat:terminal-traces} imply strong local \(L^2\)
convergence of \(\divg(\rho_tY_t)\) to
\(\divg(fY_1)\), and of the gradients of \(\theta_t\) to
that of \(\phi_1^c+c\log f\). The bounded density and field
factors converge uniformly on the compact support. Both
products in \eqref{eq:weak-hessian} therefore converge in
\(L^1\). Let \(T\uparrow1\), with \(c\) still fixed.
This proves \eqref{heat:viscous-index-estimate}. The remaining
limits will be taken successively: first \(c\downarrow0\),
and then removal of the positive terminal perturbation
\(\lambda\bar f\) by \(\lambda\downarrow0\).
\end{proof}

\Needspace{6\baselineskip}
\subsection{The zero-viscosity transport limit}

The viscous index estimate is now proved. It remains to identify the
limiting calibrated paths and to pass their index energy to the limit.
Because that energy is quadratic in the velocity, weak convergence
alone is insufficient. The action identity below will give the
required strong convergence in the density-weighted norm.

\begin{lemma}[The logarithmic heat limit]
\label{heat:laplace-principle}
As \(c\downarrow0\),
\begin{equation}
\label{heat:potential-limit}
\phi_t^c\longrightarrow\Phi_t=Q_tu
\quad\text{locally uniformly on }(0,1]\times M.
\end{equation}
At each fixed positive time the convergence is also weak in
\(W^{1,2}_{\rm loc}\).
\end{lemma}

\begin{proof}
Use the two-sided Gaussian bounds with exponents
\(4-\eta\) and \(4+\eta\), for any fixed \(0<\eta<1\),
and heat time \(ct\). The kernel is symmetric, so the volume
factor can be taken to be \(m(B_{\sqrt{ct}}(x))\) at the
evaluation point. These are the bounds of
\cite[Theorem~1.2]{JiangLiZhang2016}, after weakening the Ricci
bound to a negative one if needed.

Here are the two Laplace estimates with their uniformity
specified. Fix \(t\in[\alpha,1]\) and \(x\) in a compact
set. In the upper kernel bound reserve a small fixed fraction
\(\tau>0\) of its Gaussian exponent. Taking the infimum of
the other part of the exponent gives
\[
H_t^c(x)\le C_{\eta,\tau}
\exp\left[-\frac1{2c}
Q_{(4+\eta)t/(4(1-\tau))}u(x)\right].
\]
Indeed the integral of the reserved Gaussian, divided by
\(m(B_{\sqrt{ct}}(x))\), is uniformly bounded by the annular
volume-comparison estimate used in
Lemma~\ref{heat:kernel-ratio}.
For the lower kernel bound, choose a minimizer for the
corresponding modified quadratic cost and integrate over a
small ball of fixed radius \(r>0\) around it. Boundedness of
\(u\) keeps all these minimizers in a common compact set.
The Lipschitz bound for \(u\), and the bounded distance of
the minimizers from \(x\), change the exponent by at most
\(C_\alpha r\). The logarithm of the volume prefactor,
multiplied by \(c\), tends to zero uniformly: on the compact
region, Bishop--Gromov and full support bound small-ball
volumes below by a fixed multiple of the \(n\)-th power of
the radius, and above by a fixed finite constant. Thus
\[
\limsup_{c\downarrow0}\phi_t^c(x)
\le Q_{(4-\eta)t/4}u(x)+C_\alpha r,
\]
uniformly on the compact parameter set. Send first
\(c\downarrow0\), then \(r,\eta,\tau\downarrow0\).
The quadratic infima vary uniformly with their time parameter
on compact positive-time intervals, since their minimizers
have uniformly bounded distance. The two estimates give
\eqref{heat:potential-limit}. Finally the uniform spatial
Lipschitz bound gives weak local Sobolev compactness; closure
of the differential identifies its limit with \(\nabla\Phi_t\).
\end{proof}

We next justify compactness of the density curves, including
the second-moment issue. The relevant continuity-equation
criterion is \cite[Theorem~3.5]{GigliHan2015}: for a
\(W_2\)-continuous curve with bounded compression, its
Sobolev continuity functional controls its metric derivative.
We establish the required \(W_2\) continuity before applying
that criterion.

\begin{lemma}[Compactness with uniform moment tails]
\label{heat:curve-compactness}
For each \(0<\delta<1\), the probability curves
\(\mu_t^c=\rho_t^c m\) are uniformly Lipschitz in \(W_2\)
on \([\delta,1]\), with a constant independent of \(c\).
Their second moments are uniformly integrable. Consequently,
from any sequence \(c\downarrow0\) one can extract a
subsequence such that \(\mu_t^c\) converges uniformly in
\(W_2\) on every compact positive-time interval, including
time one, to a curve \(\mu_t=\nu_t m\). Moreover,
\(0\le\nu_t\le C_\delta\) on \([\delta,1]\).
\end{lemma}

\begin{proof}
Let \(L\) be the velocity bound in
\eqref{heat:uniform-gradients} and \(r(x)=d(o,x)\).
For a bounded nondecreasing Lipschitz function \(\chi\) that
is constant at large arguments, the continuity equation,
first with compact spatial tests and then after subtracting
that constant, gives for \(s\le t\le1\)
\[
\frac d{dt}\int\chi(r+L(t-s))\dd\mu_t^c
=\int\chi'(r+L(t-s))
\bigl(L+\langle\nabla r,v_c\rangle\bigr)\dd\mu_t^c\ge0.
\]
Compact Lipschitz tests are admitted by their strong Sobolev
test approximations. At time one, local convergence of the
density and preservation of its total mass justify the limit
for these tests. Hence
\begin{equation}
\label{heat:radial-domination}
\int\chi(r)\dd\mu_s^c
\le\int\chi(r+L(1-s))f\dd m.
\end{equation}
Approximate \(r^2\) and then
\(r^2\mathbf1_{\{r>R\}}\) monotonically by such functions.
The first choice proves finiteness of the moments, and the
second gives
\begin{equation}
\label{heat:uniform-moment-tails}
\sup_{c,\,s\in[\delta,1]}
\int_{\{r>R\}}r^2\dd\mu_s^c
\le\int_{\{r>R-L\}}(r+L)^2f\dd m\longrightarrow0.
\end{equation}
The same bound also gives tightness.

The explicit heat factors give continuity against compact
tests on the interior and at the terminal endpoint. Tightness
gives narrow continuity, and
\eqref{heat:uniform-moment-tails} upgrades it to \(W_2\)
continuity. The weak continuity equation extends to every
\(h\in W^{1,2}\), with functional
\[
\mathcal L_t^c(h)=\int\langle\nabla h,v_c\rangle\rho_t^c\dd m,
\qquad
|\mathcal L_t^c(h)|
\le\left(\int\rho_t^c|v_c|^2\dd m\right)^{1/2}
\left(\int\rho_t^c|\nabla h|^2\dd m\right)^{1/2}.
\]
To see this extension, approximate \(h\) in \(W^{1,2}\)
by compact tests; the density bound makes the errors uniform
in time on \([\delta,1]\).

For the exact domain in the cited continuity-equation
criterion, write \(S^2(M)\) for the Sobolev class with
square-integrable weak gradient, without requiring the
function itself to belong to \(L^2\). The same formula
defines \(\mathcal L_t^c\) on all of \(S^2(M)\), with the
displayed dual-seminorm bound. The integrated continuity
equation holds for every \(h\in L^1(m)\cap S^2(M)\): set
\(h_M=\max\{-M,\min\{h,M\}\}\). Then
\(|h_M|^2\le M|h|\), so \(h_M\in W^{1,2}\), and
\[
\nabla h_M=\mathbf1_{\{|h|<M\}}\nabla h,
\qquad
h_M\to h\text{ in }L^1,
\qquad
\nabla h_M\to\nabla h\text{ in }L^2.
\]
Locality makes the choice on the threshold level sets
irrelevant. Bounded compression gives
\[
\sup_{t\in[\delta,1]}
\left|\int(h_M-h)\dd\mu_t^c\right|
\le C_\delta\|h_M-h\|_1\longrightarrow0,
\]
while, on any fixed positive-time interval \(I\),
\[
\int_I|\mathcal L_t^c(h_M-h)|\dd t
\le L\sqrt{C_\delta}|I|
\|\nabla h_M-\nabla h\|_2\longrightarrow0.
\]
Thus the integrated identity passes for every pair of times,
with finite scalar integrals and an absolutely continuous
representative. The identical argument applies to the
limiting gradient functional with \(\nu,V\) below, since it
has the same density and velocity bounds. The
continuity-equation criterion now gives
\(|\dot\mu_t^c|\le L\) almost everywhere.
Finally properness, the uniform moment tails, and
Arzel\`a--Ascoli in \(\mathcal P_2(M)\) give the claimed
subsequence. The bound on the limiting densities follows by
testing their narrow limit against compact continuous
nonnegative functions. A diagonal choice handles intervals
with arbitrarily small positive left endpoint.
\end{proof}

\begin{lemma}[The action identity and its limiting equality]
\label{heat:action-limit}
Along a subsequence as in
Lemma~\ref{heat:curve-compactness}, there is a measurable
vector field \(V_t\), with \(|V_t|\le L\)
\(\nu_t\dd t\dd m\)-almost everywhere, such that
\(\partial_t\nu+\divg(\nu V)=0\). For every
\(0<\alpha<\beta\le1\),
\begin{align}
\label{heat:exact-action-identity}
\left[\int\phi_t^c\rho_t^c\dd m\right]_\alpha^\beta
&=\frac12\int_\alpha^\beta\!\int
\rho_t^c|\nabla\phi_t^c|^2\dd m\dd t\notag\\
&=\frac12\int_\alpha^\beta\!\int
\rho_t^c|v_c|^2\dd m\dd t+o(1).
\end{align}
The limiting curves are constant-speed Wasserstein geodesics
on these intervals, and
\begin{equation}
\label{heat:strong-weighted-velocity}
\int_\alpha^\beta\!\int
\rho_t^c|v_c-V_t|^2\dd m\dd t\longrightarrow0.
\end{equation}
For every fixed \(L^2\) tensor coefficient \(T\), compactly
supported in space-time in \([\alpha,\beta]\times M\),
\begin{equation}
\label{heat:strong-coefficient-velocity}
\int_\alpha^\beta\!\int
\rho_t^c|T(v_c-V_t)|^2\dd m\dd t\longrightarrow0.
\end{equation}
\end{lemma}

\begin{proof}
On \([\alpha,\beta]\) the fluxes
\(j_c=\rho^cv_c\) are bounded in the fixed Hilbert space
\(L^2(\dd t\dd m;TM)\), since
\(\int|j_c|^2\le C_\alpha\int\rho^c|v_c|^2\).
Extract a weak flux limit \(j\), and a weak-* density limit
\(\rho^c\rightharpoonup^*\nu\) in \(L^\infty(\dd t\dd m)\).
The latter agrees with the curve limit. Testing the convex
inequality \(|j_c|\le L\rho^c\) against bounded fields of
bounded support gives \(|j|\le L\nu\). Define
\(V=j/\nu\) on \(\{\nu>0\}\) and zero elsewhere. The
continuity equation passes to the limit. Its kinetic lower
semicontinuity follows, for example, from
\[
\int\nu|V|^2
=\sup_Z\int\bigl(2\langle j,Z\rangle-\nu|Z|^2\bigr),
\]
where the supremum is over bounded \(L^2\) fields with
bounded support.

For the first equality in
\eqref{heat:exact-action-identity}, write
\(a=\nabla\phi\) and use \(v_c=a+w_c\). The equations give
\[
\frac d{dt}\int\phi\rho\dd m
=\int\rho\left[-\tfrac12|a|^2+c\Delta\phi
+\langle a,a+w_c\rangle\right]\dd m
=\tfrac12\int\rho|a|^2\dd m,
\]
because
\(c\int\rho\Delta\phi=-\int\rho\langle w_c,a\rangle\).
One first uses spatial cutoffs. For fixed \(c\), their errors
vanish by mass integrability, boundedness of \(\phi,a\), and
\(|\nabla\rho|\le C_c\rho\). The identity extends to
\(\beta=1\) using the terminal traces and bounded gradients;
no Laplacian bound at exactly time one is needed. The
difference between the actions of \(a\) and \(v_c\) tends to
zero on bounded spatial regions by
Lemma~\ref{heat:weighted-estimates}. On their complement it
is bounded by a constant times the density mass, uniformly
small by \eqref{heat:uniform-moment-tails}. This proves the
second equality.

Put \(h=\beta-\alpha\) and
\[
A=\int\Phi_\beta\dd\mu_\beta-
\int\Phi_\alpha\dd\mu_\alpha.
\]
The bounded potentials, local uniform convergence, and
tightness identify this with the limit of the left side of
\eqref{heat:exact-action-identity}. The Hopf--Lax semigroup
property says
\(\Phi_\beta(y)-\Phi_\alpha(x)\le d^2(x,y)/(2h)\).
The continuity-equation criterion and action lower
semicontinuity therefore give the sandwich
\begin{align}
\label{heat:action-sandwich}
\frac{W_2^2(\mu_\alpha,\mu_\beta)}{h}
&\le\int_\alpha^\beta|\dot\mu_t|^2\dd t
\le\int_\alpha^\beta\!\int\nu|V|^2\dd m\dd t\notag\\
&\le\liminf_{c\downarrow0}
\int_\alpha^\beta\!\int\rho^c|v_c|^2\dd m\dd t
=2A\le\frac{W_2^2(\mu_\alpha,\mu_\beta)}{h}.
\end{align}
Thus all inequalities are equalities. Equality in the
metric action bound proves the geodesic assertion and
equality of the velocity norms.

To turn norm convergence into
\eqref{heat:strong-weighted-velocity}, expand the square.
For its cross term, test the weak flux convergence against
\(V\mathbf1_{B_R(o)}\); for its last term, test the density
convergence against \(|V|^2\mathbf1_{B_R(o)}\).
These are fixed admissible tests in the appropriate spaces.
Boundedness of the velocities and tightness let
\(R\to\infty\). The resulting limit is zero by equality
of actions in \eqref{heat:action-sandwich}.

Finally, for any \(A_0>0\),
\begin{align*}
\int\rho^c|T(v_c-V)|^2
\le A_0^2\int\rho^c|v_c-V|^2
+4L^2C_\alpha\int_{\{|T|>A_0\}}|T|^2.
\end{align*}
First let \(c\downarrow0\), then \(A_0\to\infty\).
This proves \eqref{heat:strong-coefficient-velocity}.
\end{proof}

\begin{lemma}[Common calibrated paths and their velocities]
\label{heat:common-paths}
There is a Borel prescription \(S:x\mapsto S(x)\in\Geo(M)\),
defined outside one \(m\)-null set and with \(S(x)_1=x\),
such that \(S(x)\) satisfies \eqref{heat:calibration}.
For every \(f\in\mathcal F\), the limiting density curves
are the marginals of
\(\pi_f=S_\#(fm)\). These plans have speed bounded by
\(\sqrt{2\operatorname{osc}(u)}\) and bounded compression
on each \([\delta,1]\). Their forward path velocities agree
with the pullback of the vector field \(V\) in
Lemma~\ref{heat:action-limit}; in particular there is no
conditional averaging loss. The same prescription defines
\(\pi_f\) for every nonnegative compactly supported
Lipschitz datum, with the same speed and positive-time
compression properties.
\end{lemma}

\begin{proof}
Initially normalize \(f\in\mathcal F\) to mass one.
Equality \eqref{heat:action-sandwich} on every positive-time
interval makes \(\mu_t\) a consistent constant-speed
Wasserstein geodesic. Its speed is bounded, so completeness
of \(\mathcal P_2(M)\) extends it to time zero. A dynamical
optimal plan represents this geodesic; one may use metric
superposition \cite[Proposition~2.1]{GigliHan2015} and equality
of endpoint cost and path action. Since
\(\|Q_tu-u\|_\infty\le t\Lip(u)^2/2\), the dual equalities
in \eqref{heat:action-sandwich} pass to time zero and give
\[
\int\Phi_1 f\dd m-\int u\dd\mu_0
=\tfrac12W_2^2(\mu_0,fm).
\]
The nonnegative pointwise duality gap therefore vanishes on
almost every path, proving \eqref{heat:calibration} for this
plan. Its individual path lengths are bounded by
\(\sqrt{2\operatorname{osc}(u)}\). The positive-time
compression is the limiting bound
\(\nu_t\le C_\delta\).

We specify the identification with \(V\). On a positive-time
interval the plan is a test plan. Its pullback velocity,
defined by differentiation of Sobolev functions, has norm
equal to the metric path speed
\cite[Theorem/Definition~2.33]{GigliDifferentialStructure}.
Let \(W\) be its conditional average given \((t,\gamma_t)\).
It solves the same continuity equation and its action is at
most the path action, by conditional Cauchy--Schwarz. The
kinetic lower bound makes this an equality. Both \(V\) and
\(W\) solve the same equation; apply the same lower bound to
\((V+W)/2\) and use the parallelogram identity. Equality of
their actions with the minimum forces \(V=W\). Equality in
conditional Cauchy--Schwarz then identifies the pullback
velocity itself with \(V_t(\gamma_t)\) almost everywhere.
Separability of the tangent module, available on this RCD
space, permits the conditional-average construction by a
countable family of module generators.

The prescription is independent of the positive datum.
Take two normalized calibrated plans, reverse their time
orientation, and form their equally weighted mixture. The
mixture is still calibrated, and hence is optimal from its
absolutely continuous initial measure, which is the mixture
of the two terminal densities, to its other endpoint
measure. The latter has finite second moment and may be
singular. Uniqueness of dynamical optimal transport from an
absolutely continuous measure
\cite[Theorem~1.1]{GigliRajalaSturm2016} forces this reversed
mixture to be induced by a map into whole geodesics.
Disintegration consequently forces both component plans to
choose the same path wherever their terminal densities
overlap. Positive members of \(\mathcal F\) are equivalent
to \(m\), so this agreement holds \(m\)-almost everywhere.

Fix the reference datum proportional to \(\bar f\) and
choose its resulting Borel map \(S\). The argument just
given shows that every heat limit for every positive datum
is \(S_\#(fm)\), after restoring mass. The prescribed-potential
uniqueness theorem \cite[Theorem~1.3]{GigliRajalaSturm2016}
applied to the \(c\)-concave potential \(Q_1u\), \(c=d^2/2\),
also gives uniqueness of the calibrated path for almost every
terminal point.

For a compact datum \(f_0\ge0\), the ratio \(f_0/\bar f\)
is bounded. Thus \(S_\#(f_0m)\) is boundedly dominated by
the reference plan. This proves its compression assertion
and transfers its path-velocity representative. In
particular one can choose this representative using the
reference plan once and for all. Restriction to a bounded
set of terminal points preserves all these statements and,
by the speed bound, gives bounded spatial support on
positive-time intervals.
\end{proof}

\begin{proof}[Proof of Theorem~\ref{heat:index}]
First take \(f\in\mathcal F\). On the common compact
support of the trial, Lemma~\ref{heat:laplace-principle}
and the uniform gradient bounds imply
\[
H_{\phi_1^c+c\log f}(Z,Z)(f)
\longrightarrow H_{\Phi_1}(Z,Z)(f).
\]
Indeed \eqref{eq:weak-hessian} pairs the gradient of its
potential with fixed \(L^2\) coefficients; weak local
Sobolev convergence suffices. The additional gradient
\(c\nabla\log f\) converges uniformly to zero.

For the right side of
\eqref{heat:viscous-index-estimate}, use
\eqref{heat:strong-coefficient-velocity} with the fixed
coefficient \(\nabla Y\). It replaces \(\nabla_{v_c}Y\)
by \(\nabla_VY\) strongly in the weighted \(L^2\) norm.
The remaining fixed squared expression is in \(L^1(\dd t\dd m)\)
and has compact support, so weak-* convergence of the
densities passes its integral to the limit. The curvature
term passes in the same way using boundedness of \(Y\) and
the velocities. Lemma~\ref{heat:common-paths} identifies
the limiting expression with the path integral in
\eqref{heat:index-inequality}. Let \(c\downarrow0\) in
Proposition~\ref{heat:viscous-index}.

Finally, for the desired compact datum \(f_0\), use
\(f_\lambda=f_0+\lambda\bar f\). The common prescription
gives
\(\pi_{f_\lambda}=\pi_{f_0}+\lambda\pi_{\bar f}\).
Both sides of the already established inequality are
linear in the terminal data. All terms involving
\(\bar f\) are finite, by compact support of the trial,
bounded speed, and positive-time compression. Let
\(\lambda\downarrow0\). This proves
\eqref{heat:index-inequality}, including the zero datum.
\end{proof}
